\section*{Bab \ref{ch:g7:volcanoes} --- Gunung Api}

\begin{solution}{exo:g7:volcanoes:1}
(a) \omterm{def:g7:volcanoes:vent}{kawah}; (b) \omterm{def:g7:volcanoes:magma}{magma}; (c) \omterm{def:g7:volcanoes:vent}{lubang letusan}; (d) \omterm{def:g4:volcanoes-earthquakes-first-look:eruption}{lava}; (e) \omterm{def:g7:volcanoes:magma-chamber}{dapur magma}.
\end{solution}

\begin{solution}{exo:g7:volcanoes:2}
(a) Efusif. (b) Eksplosif. (c) Efusif: air mancur \omterm{def:g4:volcanoes-earthquakes-first-look:eruption}{lava} melemparkan
tetes-tetes berpijar yang jatuh kembali lalu mengalir; \omterm{def:g7:volcanoes:magma}{magmanya} tidak
dihancurkan menjadi abu. (d) Eksplosif: itulah \omterm{def:g7:volcanoes:pyroclastic-flow}{aliran piroklastik}.
\end{solution}

\begin{solution}{exo:g7:volcanoes:3}
(a) Benar. (b) Salah: gas \omterm{def:g4:volcanoes-earthquakes-first-look:volcano}{gunung api} yang paling banyak adalah uap air;
karbon dioksida menyusul sesudahnya. (c) Salah: \omterm{def:g7:volcanoes:pyroclastic-flow}{aliran piroklastik} adalah
longsoran gas panas, abu, dan pecahan \omterm{def:g3:stones-pebbles-sand:rock}{batuan}. (d) Salah: \omterm{def:g4:volcanoes-earthquakes-first-look:volcano}{gunung-gunung api}
itu terletak di dalam \omterm{def:g7:plate-tectonics:plate}{lempeng} Pasifik, lebih dari 3200~km dari batas yang
terdekat, di atas sebuah \omterm{def:g7:volcanoes:hotspot}{titik panas}.
\end{solution}

\begin{solution}{exo:g7:volcanoes:4}
(a) Kira-kira 1000 sampai 1200~derajat. (b) Kira-kira 600 sampai
800~derajat. (c) Di dekat \omterm{def:g7:volcanoes:vent}{lubang letusan}, tempat \omterm{def:g4:volcanoes-earthquakes-first-look:eruption}{lavanya} kuning (1000 sampai
1200~derajat), bukan oranye (800 sampai 1000~derajat): \omterm{def:g4:volcanoes-earthquakes-first-look:eruption}{lava} mendingin
selagi mengalir menjauhi \omterm{def:g7:volcanoes:vent}{lubang letusan}.
\end{solution}

\begin{solution}{exo:g7:volcanoes:5}
(a) $838 + 221 = 1059$ \omterm{def:g4:volcanoes-earthquakes-first-look:volcano}{gunung api}. (b) $1059 \div 1214 \approx 0.872$:
kira-kira \textbf{87~\%}. (c) $69 + 18 + 12 + 0.5 = 99.5$, kira-kira
100~\%: selisih kecil itu berasal dari pembulatan setiap persentase.
\end{solution}

\begin{solution}{exo:g7:volcanoes:6}
(a) $564 \div 660 \approx 0.855$: kira-kira \textbf{85~\%}. (b)
$660 - 564 - 62 - 33 = 1$: satu kerucut curam. (c) Kerucut curam dibangun
oleh \omterm{def:g4:volcanoes-earthquakes-first-look:eruption}{lava} yang kental dan abu, ciri \omterm{def:g7:volcanoes:eruption-types}{erupsi eksplosif}: di sepanjang Cincin
Api, tempat satu \omterm{def:g7:plate-tectonics:plate}{lempeng} menunjam ke bawah \omterm{def:g7:plate-tectonics:plate}{lempeng} lain, kita dapat
mengharapkan banyak \omterm{def:g7:volcanoes:eruption-types}{erupsi eksplosif}, seperti \omterm{def:g4:volcanoes-earthquakes-first-look:eruption}{letusan} Gunung St.~Helens dan
Pinatubo.
\end{solution}

\begin{solution}{exo:g7:volcanoes:7}
(a) $1170 - 1140 = 30$~derajat. (b) $30 \div 10 = 3$~derajat per
kilometer. (c) Di dalam terowongan, \omterm{def:g4:volcanoes-earthquakes-first-look:eruption}{lava} tidak bersentuhan dengan udara:
atap yang sudah mengeras di atasnya menahan panasnya, seperti tutup panci.
Di udara terbuka, permukaannya langsung kehilangan panas ke udara.
\end{solution}

\begin{solution}{exo:g7:volcanoes:8}
(a) Kauai, paling sedikit $5.5 - 0.7 = 4.8$~juta tahun lebih tua.
(b) Kauai: \omterm{def:g4:clouds-rain-snow:precipitation}{hujan} dan ombak sudah mengikisnya jauh lebih lama. (c) Pulau
Besar, tempat \omterm{def:g4:volcanoes-earthquakes-first-look:eruption}{lava} baru masih keluar. (d) Ke arah barat laut: Kauai, yang
lebih dahulu dibawa menjauhi \omterm{def:g7:volcanoes:hotspot}{titik panas}, kini berada di sebelah barat
\omterm{def:g2:water-on-earth:ocean}{lautnya}.
\end{solution}

\begin{solution}{exo:g7:volcanoes:9}
(a) Hampir semua \omterm{def:g4:volcanoes-earthquakes-first-look:volcano}{gunung api} terletak di \omterm{def:g7:plate-tectonics:plate-boundary}{batas lempeng} (87~\% dalam daftar
itu): \omterm{def:g4:volcanoes-earthquakes-first-look:volcano}{gunung api} yang begitu jauh dari batas mana pun adalah pengecualian.
(b) Jauh di dalam \omterm{def:g6:inside-the-earth:layers}{mantel} di bawah Hawaii, sebuah \omterm{def:g7:volcanoes:hotspot}{titik panas} tetap panas
selama jutaan tahun. \omterm{def:g7:volcanoes:hotspot}{Titik panas} itu melelehkan sebagian \omterm{def:g3:stones-pebbles-sand:rock}{batuan} di atasnya.
\omterm{def:g7:volcanoes:magma}{Magmanya}, yang lebih ringan daripada \omterm{def:g3:stones-pebbles-sand:rock}{batuan} padat di sekelilingnya, naik
menembus \omterm{def:g7:plate-tectonics:plate}{lempeng} dan meletus. Selagi \omterm{def:g7:plate-tectonics:plate}{lempeng} bergerak, \omterm{def:g4:volcanoes-earthquakes-first-look:volcano}{gunung-gunung api}
baru tumbuh di atas \omterm{def:g7:volcanoes:hotspot}{titik panas} dan \omterm{def:g4:volcanoes-earthquakes-first-look:volcano}{gunung-gunung api} yang lama terbawa
pergi. (c) Titik-titik oranye (``di dalam \omterm{def:g7:plate-tectonics:plate}{lempeng}'').
\end{solution}

\begin{solution}{exo:g7:volcanoes:10}
\omterm{def:g7:volcanoes:magma}{Magma} yang pertama sudah kehilangan sebagian besar gasnya dalam
perjalanan ke atas, seperti sebotol air soda yang dibiarkan terbuka atau
dibiarkan berdiri: ketika \omterm{def:g7:volcanoes:magma}{magma} itu mencapai permukaan, tidak ada lagi
yang dapat meledakkannya, sehingga \omterm{def:g7:volcanoes:magma}{magma} itu merembes keluar. \omterm{def:g7:volcanoes:magma}{Magma} tanggal
12 dan 15~Juni masih menyimpan gasnya, seperti botol yang dikocok: ketika
tekanannya turun di dekat permukaan, gas itu serentak membentuk
gelembung-gelembung, yang membesar dan mengoyak \omterm{def:g7:volcanoes:magma}{magma} yang kental itu
menjadi abu. \omterm{def:g7:volcanoes:magma}{Magma} yang pertamalah yang sudah kehilangan gasnya.
\end{solution}

\begin{solution}{exo:g7:volcanoes:11}
(a) $10 \div 80 = 0.125$~jam, dan $0.125 \times 60 = 7.5$~menit. (b)
$10 \div 130 \approx 0.077$~jam, kira-kira $0.077 \times 60 \approx
4.6$~menit, jadi kira-kira 5~menit. (c) Keluarga itu memerlukan 20~menit,
tetapi aliran itu tiba dalam 5 sampai 8~menit: begitu aliran itu terlihat,
sudah terlambat. Keluarga itu hanya dapat selamat jika mereka pergi ketika
para ilmuwan melihat tanda-tanda peringatan, sebelum \omterm{def:g4:volcanoes-earthquakes-first-look:eruption}{letusan}.
\end{solution}

\begin{solution}{exo:g7:volcanoes:12}
(a) $111 \times 360 = 39\,960$: kira-kira \textbf{40\,000~km}. (b)
$40\,000 \div 15 \approx 2667$: kira-kira 2\,670~km per hari. (c)
$2667 \div 24 \approx 111$: kira-kira \textbf{110~km per jam}. (d) Oleh
angin: abu yang halus tidak mempunyai kecepatan sendiri setelah \omterm{def:g4:volcanoes-earthquakes-first-look:eruption}{letusan}
melemparkannya ke atas; abu itu melayang tinggi di udara dan angin
membawanya, sementara abu itu perlahan-lahan jatuh.
\end{solution}

\begin{solution}{pb:g7:volcanoes:1}
\textbf{1.} Dari 16 sampai 31~Maret: $31 - 16 = 15$~hari; lalu 30~hari di
bulan April dan 18~hari di bulan Mei: $15 + 30 + 18 = 63$~hari.

\textbf{2.} $10\,000 \div 62 \approx 161$: lebih dari 160 \omterm{def:g4:volcanoes-earthquakes-first-look:earthquake}{gempa bumi} kecil
sehari, rata-rata.

\textbf{3.} Penggembungan itu menunjukkan bahwa \omterm{def:g7:volcanoes:magma}{magma} sudah naik tinggi di
dalam \omterm{def:g1:landscapes-around-us:landform}{gunung} dan mendorong \omterm{def:g3:stones-pebbles-sand:rock}{batuan} ke luar. Penggembungan itu juga membuat
lereng utara curam dan tidak stabil, siap untuk longsor.

\textbf{4.} Segitiga biru terletak di tempat satu \omterm{def:g7:plate-tectonics:plate}{lempeng} menunjam ke bawah
\omterm{def:g7:plate-tectonics:plate}{lempeng} lain. Kita dapat mengharapkan \omterm{def:g7:volcanoes:magma}{magma} yang kental dan kaya gas,
\omterm{def:g7:volcanoes:eruption-types}{erupsi eksplosif}, dan kerucut yang curam (\omterm{def:g4:volcanoes-earthquakes-first-look:volcano}{gunung api} strato) --- tepat
seperti yang ditunjukkan foto pertama.

\textbf{5.} $240 \div 60 = 4$~kilometer per menit.

\textbf{6.} Lereng yang menggembung itu menahan \omterm{def:g7:volcanoes:magma}{magma} di dalam, di bawah
tekanan, seperti tutup botol. Ketika lereng itu longsor, tekanannya
seketika turun; gas yang terlarut di dalam \omterm{def:g7:volcanoes:magma}{magma} serentak berubah menjadi
gelembung-gelembung dan menyembur keluar melalui sisi yang terbuka ---
ke samping, tempat tutupnya hilang.

\textbf{7.} $27 \div 480 \approx 0.056$~jam, dan $0.056 \times 60 \approx
3.4$~menit. Karena ledakan itu bergerak paling sedikit secepat itu,
waktunya paling lama kira-kira 3.4~menit.

\textbf{8.} $8 \div 130 \approx 0.062$~jam, dan $0.062 \times 60 \approx
3.7$: kira-kira \textbf{4~menit}.

\textbf{9.} $24 \div 8.8 \approx 2.7$: tiang abu itu naik kira-kira 2.7 kali
lebih tinggi daripada Gunung Everest.

\textbf{10.} Efusif: \omterm{def:g4:volcanoes-earthquakes-first-look:eruption}{lava} yang kental itu keluar tanpa meledak dan menumpuk
di atas \omterm{def:g7:volcanoes:vent}{lubang letusan}. Seperti botol yang dibiarkan berdiri, \omterm{def:g7:volcanoes:magma}{magma} ini
sudah kehilangan sebagian besar gasnya.

\textbf{11.} Sebelumnya, \omterm{def:g1:landscapes-around-us:landform}{gunung} itu adalah kerucut runcing yang mulus dan
tertutup \omterm{def:g4:clouds-rain-snow:precipitation}{salju}, di atas hutan-hutan hijau. Sesudahnya, puncak dan lereng
utaranya hilang: \omterm{def:g7:volcanoes:vent}{kawah} yang sangat besar, yang terbuka ke arah kita,
menggantikannya, dan lereng-lereng serta \omterm{def:g1:landscapes-around-us:landform}{lembah} di depannya kelabu oleh abu
dan pecahan \omterm{def:g3:stones-pebbles-sand:rock}{batuan}.

\textbf{12.} $2950 - 2549 = 401$: \omterm{def:g1:landscapes-around-us:landform}{gunung} itu kehilangan \omterm{def:g4:world-physical-map:altitude}{ketinggian}
kira-kira \textbf{401~m}, kira-kira 400~meter, seperti yang diceritakan di
kelas~4.
\end{solution}
